\documentclass[11pt]{article}
\usepackage[a4paper,margin=28mm]{geometry}
\usepackage{amsmath,amssymb,amsthm,booktabs}
\usepackage[hidelinks]{hyperref}
\newtheorem{theorem}{Theorem}
\newtheorem{lemma}[theorem]{Lemma}
\newtheorem{corollary}[theorem]{Corollary}
\newtheorem{proposition}[theorem]{Proposition}
\theoremstyle{remark}
\newtheorem{remark}[theorem]{Remark}
\newcommand{\K}{\mathcal K}
\newcommand{\A}{\mathcal A}
\newcommand{\NN}{\mathbb Z_{\geq0}}
\newcommand{\RR}{\mathbb R}
\newcommand{\bd}{\partial_{\min}}
\newcommand{\affdim}{\operatorname{affdim}}
\newcommand{\codim}{\operatorname{codim}}
\title{Local separation and M\"obius inversion for pyramidal partitions}
\author{}
\date{17 October 2026}
\begin{document}
\maketitle
\begin{center}
\small\textbf{Research draft: priority audit incomplete; not for submission.}
\end{center}

\begin{abstract}
A pyramidal partition is a finite subset of the nonnegative integer lattice
separated from its complement by a hyperplane. We give an infinite family of
three-dimensional plane partitions whose coordinate-plane sections are all
triangular although the full partition is not pyramidal. An equal-sum pair
obstruction disproves the slice criterion stated in Galv\'an's 2023 thesis.
We also represent every containment
interval of these partitions as a lower interval in the region poset of a finite
central hyperplane arrangement. Applying Edelman's M\"obius-function theorem
gives an explicit formula in every dimension: a value is nonzero exactly when
the interval adds one complete lattice layer on a positive cutting hyperplane;
its sign is determined by the affine dimension of that layer. The criterion
reduces evaluation to rational linear feasibility and rank, without enumerating
the interval. As a consequence, the number of pyramidal subpartitions satisfies
an alternating recurrence over removable exposed faces; in dimension three
only linearly many terms can contribute at each top partition.
Exact checks on 27,100 finite intervals support the implementation.
The arrangement theorem is established prior work; the precise application and
its priority require the audit stated at the end of this draft.
\end{abstract}

\section{Two questions about positive lattice cuts}

Fix $k\geq1$. A \emph{corner cut}, or \emph{pyramidal partition}, is a finite set
\[
 D(w,t)=\{a\in\NN^k:w\cdot a<t\},\qquad w\in\RR_{>0}^k,
\]
where the cutting hyperplane contains no lattice point. We include the empty
cut and order the resulting set $\K_k$ by inclusion. A weak cutoff gives the
same collection of finite sets, by a small change of threshold. Translation
by $(1,\ldots,1)$ identifies this convention with positive-coordinate
definitions. Corner cuts were studied by Onn and Sturmfels~\cite{OS}.

In dimension two these sets are triangular partitions. Every coordinate-plane
section of a three-dimensional corner cut is triangular or empty, since fixing
one coordinate in its defining inequality leaves positive weights on the other
two. Conjecture~4.1 of Galv\'an's thesis~\cite{GalvanThesis}, p.~51, asks
whether this necessary condition is sufficient for a plane partition.
Theorem~\ref{thm:slices} gives a uniform counterexample. Thus separate
two-dimensional cutting lines need not admit a common three-dimensional plane.
The obstruction uses elementary convex separation; we do not claim a new
separation criterion.

Every interval $[A,B]$ is finite, being a collection of subsets of $B$.
Its incidence M\"obius function is defined by
\[
 \mu(A,A)=1,\qquad
 \mu(A,B)=-\sum_{A\subseteq D\subsetneq B}\mu(A,D)\quad(A\subsetneq B),
\]
with $D$ ranging over corner cuts. This is not the arithmetic M\"obius function.
Section~8.2 of the Elizalde--Galv\'an preprint~\cite{EG} reports experimental
evidence for values in $\{-1,0,1\}$ in dimension three. The following formula
addresses that observation in all dimensions by applying an older theorem.

\begin{theorem}\label{thm:main}
Let $A\subsetneq B$ in $\K_k$ and $U=B\setminus A$. Then $\mu(A,B)\ne0$
if and only if there exist $w\in\RR_{>0}^k$ and $t\in\RR$ such that
\begin{equation}\label{eq:layer}
 \begin{cases}
 w\cdot a<t & (a\in A),\\
 w\cdot u=t & (u\in U),\\
 w\cdot b>t & (b\in\NN^k\setminus B).
 \end{cases}
\end{equation}
When this holds,
\begin{equation}\label{eq:mu}
 \mu(A,B)=(-1)^{1+\affdim(U)}.
\end{equation}
In particular, all values belong to $\{-1,0,1\}$ for every $k$.
\end{theorem}

The condition says that $A$ consists of the points strictly below one positive
cutting hyperplane and $B$ additionally contains its entire lattice layer.
No lattice assumption on the containment poset is needed.

\section{Triangular sections do not imply a pyramidal partition}

For a finite downset $D\subset\NN^3$, a coordinate-plane section is obtained
by fixing one coordinate and retaining the other two in their original order.

\begin{theorem}\label{thm:slices}
For every integer $m\geq1$ there is a downset $D_m\subset\NN^3$ of
cardinality $4m+5$ such that every coordinate-plane section is a corner cut,
but $D_m$ is not separated from its complement by any affine hyperplane.
The same failure of recognition by coordinate two-dimensional sections occurs
in every dimension $k\geq3$.
\end{theorem}
\begin{proof}
Define $D_m$ by its $z$-height matrix, with columns indexed by $x=0,1$
and rows by $y=0,1,2$:
\[
 H_m=\begin{pmatrix}2m+2&m\\m+1&1\\1&0\end{pmatrix},\qquad
 D_m=\{(x,y,z):0\leq z<H_m(y,x)\}.
\]
All other heights are zero. The heights decrease in each coordinate, so this
is a downset; their sum is $4m+5$.

Table~\ref{tab:slices} gives positive integer weights $c$ and twice a threshold
$2t$ for every nonempty section. In each row the section is exactly
$\{v\in\NN^2:c\cdot v<t\}$. For example, at $x=0$ the inequality
$(2m+1)y+2z<(8m+5)/2$ allows $2m+2,m+1,1$ values of $z$ at
$y=0,1,2$, respectively, and no values at larger $y$.
The other rows follow by the same integer inequalities, or directly from
the listed heights. Every twice-threshold is odd, so no lattice point
lies on a cutting line. All omitted sections are empty and use $c=(1,1)$,
$t=-1/2$.

\begin{table}[ht]
\centering
\begin{tabular}{lll}
\toprule
fixed coordinate & remaining weights $c$ & $2t$\\
\midrule
$x=0$ & $(2m+1,2)$ & $8m+5$\\
$x=1$ & $(2m-1,2)$ & $4m-1$\\
$y=0$ & $(2m+3,2)$ & $8m+5$\\
$y=1$ & $(m,1)$ & $2m+1$\\
$y=2$ & $(1,1)$ & $1$\\
$z=0$ & $(3,2)$ & $11$\\
$1\leq z<m$ & $(1,1)$ & $3$\\
$z=m$ & $(2,1)$ & $3$\\
$m<z\leq2m+1$ & $(1,1)$ & $1$\\
\bottomrule
\end{tabular}
\caption{Exact cutting certificates for the sections of $D_m$.
Remaining coordinates are in their original order.}\label{tab:slices}
\end{table}

On the other hand, the points
\[
 a=(0,0,2m+1),\quad d=(1,1,0)
\]
belong to $D_m$, whereas
\[
 b=(1,0,m),\quad c=(0,1,m+1)
\]
do not. They satisfy $a+d=b+c$. If a strict affine cutting inequality
contained $a,d$ and excluded $b,c$, summing it over each pair would give
contradictory inequalities for the same vector sum. Equivalently their two
segments meet at $(1/2,1/2,m+1/2)$. Hence even a separator without a
positivity restriction is impossible.

For $k>3$, take $D_m\times\{0\}^{k-3}$. The same four points rule out
separation. A coordinate two-dimensional section is empty, a section from the
three-dimensional example, a finite initial segment on one coordinate axis,
or a singleton. Each is a two-dimensional corner cut: for an axis segment,
choose the other weight sufficiently large. This proves the extension.
\end{proof}

For $m=1$ the example has nine points and height matrix
$\left(\begin{smallmatrix}4&1\\2&1\\1&0\end{smallmatrix}\right)$.
We make no minimality claim. The conclusion directly contradicts
Conjecture~4.1 as stated in the thesis; priority relative to later work
remains subject to the audit below. It does not say that every proper
higher-dimensional section of the embedded examples is pyramidal.

\section{A finite arrangement for each interval}

For a finite downset $B\subset\NN^k$, let
\[
 E=\bd B=\min(\NN^k\setminus B)
\]
denote the minimal excluded border, with respect to coordinatewise order.
Every point outside $B$ dominates an element of $E$. The border is finite:
if $B\ne\varnothing$, every border point is $b+e_i$ for some $b\in B$.
For the empty set, the border consists of the origin.

Consider the following finite central arrangement in $V=\RR^{k+1}$, with
coordinates $(w,t)$:
\begin{equation}\label{eq:arrangement}
 \A_B=\{w_i=0:1\leq i\leq k\}
 \ \cup\ \{h_a=0:a\in B\cup E\},\qquad h_a(w,t)=w\cdot a-t.
\end{equation}
It is essential, since the coordinate normals and any normal $(a,-1)$ span $V$.
A region is a connected component of the complement of the hyperplanes.
Its strict sign vector determines it uniquely, since each sign cell is convex.

For a base region $R_0$, let $S(R_0,R)$ be the set of hyperplanes separating
$R_0$ and $R$. The \emph{region poset} $P(\A_B,R_0)$ orders regions by
inclusion of these separating sets.

\begin{lemma}\label{lem:interval}
For $A\subseteq B$ in $\K_k$, the interval $[A,B]$ is isomorphic to a lower
interval $[R_A,R_B]$ in $P(\A_B,R_A)$.
\end{lemma}

\begin{proof}
Strict cutting parameters for $A$ and $B$ define regions $R_A$ and $R_B$.
All coordinate signs are positive. For a point $a\in B\cup E$, the sign
of $h_a$ is negative precisely when $a$ belongs to the cut. Hence
\[
 S(R_A,R_B)=\{h_u=0:u\in B\setminus A\}.
\]
A region $R$ in the indicated lower interval therefore has $w_i>0$, has
$h_a<0$ for $a\in A$, and has $h_b>0$ for $b\in E$. Let $(w,t)\in R$.
The set $D(w,t)$ is finite and contains $A$. If $b\notin B$, choose
$e\in E$ with $e\leq b$. Positivity gives
$w\cdot b-t\geq w\cdot e-t>0$, so $D(w,t)\subseteq B$.
These inequalities also show that its cutting hyperplane has no lattice
point: points of $B$ have specified nonzero signs, and all other points
have positive signs.

Conversely, every cut $D$ with $A\subseteq D\subseteq B$ gives a strict
sign vector for \eqref{eq:arrangement}, and thus exactly one region in
$[R_A,R_B]$. Its separating set consists of the hyperplanes indexed by
$D\setminus A$. The two maps are inverse and preserve inclusion.
\end{proof}

\section{Edelman's facial criterion}

We use the following established result, in the notation above.
Edelman's Theorem~1.11~\cite{Edelman} states that $\mu(R_0,R)$ is zero
unless $S(R_0,R)$ is exactly the set of arrangement hyperplanes containing
a nonempty face $F$ of $\overline{R_0}$. In that case its value is
$(-1)^{\codim F}$. Here a face includes the closed region itself.
The result applies to the essential arrangement \eqref{eq:arrangement}.

\begin{proof}[Proof of Theorem~\ref{thm:main}]
Use Lemma~\ref{lem:interval}, with $U=B\setminus A$.
If the M\"obius value is nonzero, Edelman's criterion gives a face $F$
whose containing hyperplanes are exactly $\{h_u=0:u\in U\}$.
At a point in $\operatorname{relint}F$, every other arrangement sign
remains the strict sign of $R_A$. In particular $w_i>0$, $h_a<0$ for
$a\in A$, and $h_b>0$ for $b\in E$. The border argument in the lemma
extends the last inequality to all points outside $B$, yielding
\eqref{eq:layer}.

Conversely, suppose $(w,t)$ satisfies \eqref{eq:layer}. It has zeros
exactly on the arrangement hyperplanes indexed by $U$, and the strict
$R_A$ signs elsewhere. It lies in $\overline{R_A}$: adding $\epsilon$
times a point of $R_A$ makes all the zero signs agree with $R_A$, while
the other signs persist for sufficiently small $\epsilon>0$.
Consequently it belongs to the relative interior of a face $F$ of
$\overline{R_A}$ contained in exactly the hyperplanes indexed by $U$.
Edelman's criterion applies.

The normals of these hyperplanes are $(u,-1)$, $u\in U$. Their rank is
$1+\affdim(U)$: subtract one fixed row $(u_0,-1)$ from the other rows.
This rank is $\codim F$, proving \eqref{eq:mu}. The face has dimension
at least one, since it contains a ray with strictly positive $w$.
\end{proof}

\begin{corollary}
Let $B_m=\{a\in\NN^k:\sum_i a_i\leq m\}$ for $m\geq0$. For $m\geq1$,
\[
 \mu(B_{m-1},B_m)=(-1)^k.
\]
\end{corollary}
\begin{proof}
Take $w=(1,\ldots,1)$ and $t=m$. The layer of degree $m$ has affine
dimension $k-1$, including when $k=1$.
\end{proof}

For example, in dimension three the interval $[B_1,B_2]$ adds six points
and has value $-1$, determined by their two-dimensional affine span.
The exact verifier finds 26 cuts in this interval; it is not a Boolean
interval on six independently selectable points.

The word ``entire'' in the layer criterion is necessary. In dimension two,
take $A=B_1$ and $B=A\cup\{(2,0),(1,1)\}$. Both are corner cuts, but
their difference is only part of the line $x+y=2$: the excluded point
$(0,2)$ lies on the same line. Equality on the two added points forces
$w_1=w_2$, violating the strict inequality at $(0,2)$. Thus $\mu(A,B)=0$.
Indeed the interval is the three-element chain obtained by adding $(2,0)$
before $(1,1)$. Adding the entire degree-two line instead gives value $+1$.

\begin{proposition}\label{prop:fouratoms}
There is a three-dimensional interval with four atoms and M\"obius value
$-1$ whose top is not the join of any two or three atoms. Both endpoints
are exponent prefixes for the prime set $\{3,5,13\}$.
\end{proposition}
\begin{proof}
Let $A=\{(x,y,z)\in\NN^3:2x+3y+5z<11\}$ and
$B=\{(x,y,z)\in\NN^3:2x+3y+5z\leq11\}$. The layer $B\setminus A$ is
\[
 (0,2,1),\quad(1,3,0),\quad(3,0,1),\quad(4,1,0).
\]
These points form a parallelogram, with the sum of its first and fourth
points equal to the sum of its second and third points. A small perturbation
of the cutting hyperplane realizes exactly the strictly linearly separable
subsets of these four points. All subsets occur except the two opposite
pairs. Indeed adjacent pairs and all singletons and triples are separated
by lines in the plane of the layer, while opposite pairs have intersecting
convex hulls. A sufficiently small perturbation preserves every strict
inequality off the layer, as can be checked on the finite borders.

Thus the interval has 14 elements and four atoms. Every three atoms have
their three-point subset as a proper upper bound, so the top cannot be
the join of two or three atoms. The M\"obius recurrence gives $-1$ at
each singleton, $+1$ at each of the four adjacent pairs, zero at each
triple, and $-1$ at the top. This agrees with Theorem~\ref{thm:main}.

Finally, $|A|=20$ and $|B|=24$. Exact increasing smooth enumeration for
$\{3,5,13\}$ gives the values at ranks 20 through 25 as
\[
 243,\ 325,\ 351,\ 375,\ 405,\ 507.
\]
The first 20 exponent vectors are $A$, and the next four are the layer
displayed above. Independently of heap enumeration, direct evaluation on
the maximal elements of $A$ gives maximum product $243$, and evaluation
on the minimal excluded border of $B$ gives minimum product $507$.
The four layer products lie strictly between them. Monotonicity of
positive prime powers extends these border inequalities to all included
and excluded vectors. The verifier saves both finite border certificates,
establishing the claimed realization.
\end{proof}

This also contradicts the value rule in Conjecture~4.5 of
Galv\'an's thesis~\cite{GalvanThesis}, p.~60: that rule assigns zero
when the interval is neither a cover nor has its top generated by two
or three atoms. The displayed top has value $-1$ and satisfies none of
those conditions. The thesis also conjectures that the full poset is a
lattice; its failure was already reported in the later preprint~\cite{EG}.
We claim no new discovery of that failure.

\section{A recurrence over removable exposed faces}

For $B\in\K_k$, put $I(B)=|\{D\in\K_k:D\subseteq B\}|$, including
the empty cut, so $I(\varnothing)=1$. For nonempty $B$, let $P=\operatorname{conv}B$.
Call a nonempty face $F$ of $P$ \emph{admissible} if some $w>0$ satisfies
\begin{equation}\label{eq:admissible}
 F=\{x\in P:w\cdot x=t\},\qquad
 t=\max_{a\in B}w\cdot a<\min_{e\in\bd B}w\cdot e.
\end{equation}
Here the whole polytope is permitted as a face. Write $\mathcal F(B)$ for
these faces and $U_F=B\cap F$ for their lattice layers.

\begin{corollary}\label{cor:recurrence}
For every nonempty $B\in\K_k$,
\begin{equation}\label{eq:recurrence}
 I(B)=1+\sum_{F\in\mathcal F(B)}(-1)^{\dim F}I(B\setminus U_F).
\end{equation}
The proper nonzero entries $\mu(A,B)$ correspond bijectively to
$\mathcal F(B)$. If $n=|B|$, their number is at most
$\sum_{j=1}^{\min(k,n)}\binom{n}{j}$. For $k=3$ it is at most $6n$;
if $P$ is three-dimensional with $v$ vertices, the sharper bound is $6v-10$.
\end{corollary}
\begin{proof}
Theorem~\ref{thm:main} identifies a nonzero proper interval with a layer
$U=B\setminus A$ maximizing a positive linear functional on $B$ and strictly
below its excluded border. Its convex hull is an admissible exposed face,
and $U=B\cap\operatorname{conv}U$. Conversely an admissible face satisfies
the theorem's layer criterion. A threshold just below its score shows that
$B\setminus U_F$ is a corner cut. The two constructions are inverse.

M\"obius inversion gives
\[
 \sum_{A\subseteq B}\mu(A,B)I(A)
 =\sum_{D\subseteq B}\ \sum_{D\subseteq A\subseteq B}\mu(A,B)=1,
\]
where every summation is over corner cuts. Isolating $A=B$ and substituting
the layer formula gives \eqref{eq:recurrence}.

Every admissible face has affine dimension at most $k-1$, and has an affine
basis of at most $k$ points of $B$. It is determined by that basis: for any
face $F$ of a polytope, $P\cap\operatorname{aff}F=F$. Selecting one basis
per face therefore gives the stated binomial bound. In three dimensions,
Euler's relation and $e\leq3v-6$ give $v+e+f=2e+2\leq6v-10$ nonempty
proper faces of a full-dimensional polytope. If $|B|>1$, the whole $P$
is not admissible: $B$ contains the origin and a nonzero nonnegative vector,
on which a positive functional has different values. If $P$ has dimension
two or one, there are at most $2v$ or two nonempty proper faces, respectively.
For a singleton there is one admissible face. These cases give $6n$ uniformly.
\end{proof}

This is standard M\"obius inversion applied to the layer formula, not a new
inversion theorem. It extends the subpartition recurrence of
Elizalde--Galv\'an~\cite[Lemma~7.1 in the preprint]{EG} beyond dimension two,
using all admissible faces rather than a bound on the number of removable
cells. The bound concerns terms at one top partition; it does not bound the
number of distinct smaller cuts reached recursively or establish a faster
global enumeration algorithm.

For fixed $k$, all terms can be found from an explicit list of $B$ by
enumerating affine bases of size at most $k$ among its coordinatewise maximal
points. For each basis $S$, form $U=B\cap\operatorname{aff}S$, remove
duplicates, and test the strict layer system below for $(B\setminus U,B)$.
Every admissible layer occurs, because positivity makes its points
coordinatewise maximal. Exact affine rank and rational linear programming
give a polynomial-time procedure in the explicit input for fixed dimension.
The supplied verifier uses Fourier--Motzkin elimination and makes no such
claim in variable dimension.

For example, let $B=\{(x,y,z)\in\NN^3:x+y+z\leq2\}$. Its admissible
faces are the three vertices, three edges and triangular face on $x+y+z=2$.
The contained-cut counts after deleting these layers are respectively
$32,18,9$, constant within each face dimension. Thus
\[
 I(B)=1+3\cdot32-3\cdot18+9=52.
\]
An edge layer here has three lattice points, and the face layer has six;
removing only the vertices of those faces would give different sets.

\section{Evaluation and reproducible checks}

The formula does not require constructing the containment interval.
Given explicit point lists for $A\subsetneq B$, fix $u_0\in B\setminus A$.
Condition \eqref{eq:layer} is equivalent to
\begin{align*}
 w_i&>0 &&(1\leq i\leq k),\\
 (u_0-a)\cdot w&>0 &&(a\in\max A),\\
 (b-u_0)\cdot w&>0 &&(b\in\bd B),\\
 (u-u_0)\cdot w&=0 &&(u\in B\setminus A).
\end{align*}
Positivity extends the inequalities from the two borders to the full
sets. Since all strict inequalities are homogeneous and finite in number,
one may replace their right sides by $1$ after a common positive scaling.
This gives a rational linear program in $k$ variables. A rational rank
calculation determines the sign. The border has at most $k|B|$ candidates,
so both computations have polynomial encoding size in the explicit input.
This input convention does not describe a large prefix merely by a binary rank.

The dependency-free Python artifact
\texttt{scripts/run\_pyramidal\_poset\_checks.py} enumerates every downset
inside each top in Table~\ref{tab:checks}. It recognizes corner cuts by
exact strict separation and computes every comparable interval value by
the defining M\"obius recurrence. It compares these values with the
cutting-layer criterion. Feasibility uses rational Fourier--Motzkin
elimination, with integer witnesses verified by substitution. Independent
vertex enumeration cross-checks 233 nonempty downsets in the smaller
fixtures. This implementation is a small-case verifier; its running time
is not claimed polynomial in dimension.

\begin{table}[ht]
\centering
\begin{tabular}{rrrrr}
\toprule
$k$ & top degree & downsets & corner cuts & intervals\\
\midrule
1&6&8&8&36\\
2&4&132&48&797\\
3&2&96&52&736\\
3&3&2498&233&11799\\
4&2&1338&300&13732\\
\bottomrule
\end{tabular}
\caption{Exhaustive checks inside the tops $B_m$. All 27,100 comparable
intervals agree with the formula. This is not an enumeration of all cuts
of a prescribed cardinality.}\label{tab:checks}
\end{table}

Run from the repository root:
\begin{quote}\ttfamily\small
python3 -B scripts/run\_pyramidal\_poset\_checks.py
\end{quote}
The saved JSON file \texttt{results/research/pyramidal\_poset.json}
contains the exact top points, masks of every recognized cut, all interval
values, representative cutting witnesses, command, Python version and
source hashes. The checks support the implementation; the uniform result
is the proof above.

The same command verifies the slice family for $1\leq m\leq32$.
It checks $1,280$ nonempty coordinate sections, their supplied integer
certificates against included points and the entire minimal excluded border,
and independently computed rational cutting witnesses. It also verifies
global infeasibility and the equal-sum obstruction for every parameter.
The saved file contains all point sets and certificates. These finite checks
verify the implementation; Table~\ref{tab:slices} and the balanced pairs
prove Theorem~\ref{thm:slices} for every $m$.

For every recognized top in the five fixtures, the verifier also enumerates
candidate removable faces from affine bases, independently of the list of
subcuts. It compares the recovered terms with the entire nonzero M\"obius
column, and evaluates \eqref{eq:recurrence} against the direct number of
contained cuts. Exact face weights, affine bases, coefficients and counts
are saved with the report. Thus omission of a contributing face, as well
as an incorrect count or sign, causes a check to fail.

\section{Connection to the originating sequence problem}

For distinct primes $p_1,\ldots,p_k$, write the increasing smooth numbers
as $p_1^{a_1}\cdots p_k^{a_k}$. Their first $N$ exponent vectors form
a corner cut for $w_i=\log p_i$, with a threshold between consecutive
scores. Thus their prefix sets belong to $\K_k$. The theorem concerns
the containment structure across such weighted exponent cuts, rather than
only the chain obtained from one fixed prime set.

The argument uses positivity and linear ordering, not a further arithmetic
property of primes. It proves no improvement in locating primes and no
new random-access complexity bound for smooth numbers. Its intended use
is the geometric description and evaluation of intervals of pyramidal
partitions, an existing combinatorial question.

\section*{Priority and readiness audit}

\begin{quote}
\small\textbf{Working assumption (17 October 2026).}
At the user's request, ongoing work assumes that the relevant results in the
2025 journal article agree with the accessible 2023 v1 preprint. This is a
provisional assumption, not a verified comparison. The paywall does not halt
research under this assumption. Before submission, compare the
\href{https://doi.org/10.1016/j.aam.2024.102807}{published article}
with \href{https://arxiv.org/abs/2312.16353v1}{arXiv:2312.16353v1},
especially the M\"obius results in Section~4, Lemma~7.1, the pyramidal
discussion in Section~8.2, and any added references or acknowledgements.
\end{quote}

The candidate contributions are the slice counterexample family of
Theorem~\ref{thm:slices}, Lemma~\ref{lem:interval}, and the explicit
single-layer interpretation of Edelman's theorem for pyramidal partitions.
Edelman's region-poset formula is not a new result of this project.
The motivating observation has been checked in the 2023 preprint;
the relevant definitions, conjectures and conclusion of Galv\'an's full
2023 thesis have now been inspected. The constructions refute its stated
Conjectures~4.1 and~4.5. The full 2025 journal text and subsequent work still
require comparison before a claim of a previously unrecorded resolution is made.
General interval homotopy, a new arrangement theorem, and prime-distribution
consequences are not claimed. Authorship, disclosure of computational and
AI assistance, and journal requirements must be settled before submission.
This draft is therefore not yet submission-ready.

\begin{thebibliography}{9}
\bibitem{Edelman}
P.~H. Edelman,
A partial order on the regions of $\RR^n$ dissected by hyperplanes,
\emph{Trans. Amer. Math. Soc.} \textbf{283} (1984), 617--631.
\href{https://doi.org/10.1090/S0002-9947-1984-0737888-6}{doi:10.1090/S0002-9947-1984-0737888-6}.

\bibitem{EG}
S.~Elizalde and A.~B. Galv\'an,
Triangular partitions: enumeration, structure, and generation,
\emph{Adv. Appl. Math.} \textbf{163} (2025), 102807.
\href{https://doi.org/10.1016/j.aam.2024.102807}{doi:10.1016/j.aam.2024.102807}.
The version inspected here is
\href{https://arxiv.org/abs/2312.16353v1}{arXiv:2312.16353v1} (2023),
especially Section~8.2.

\bibitem{OS}
S.~Onn and B.~Sturmfels,
Cutting corners,
\emph{Adv. Appl. Math.} \textbf{23} (1999), 29--48.
\href{https://doi.org/10.1006/aama.1999.0645}{doi:10.1006/aama.1999.0645}.

\bibitem{GalvanThesis}
A.~B. Galv\'an P\'erez-Ilzarbe,
\emph{Properties of triangular partitions and their generalizations},
bachelor's thesis, Universitat Polit\`ecnica de Catalunya, 2023.
\href{https://upcommons.upc.edu/entities/publication/99559712-a4e1-46ae-99b3-88546b90b4f3}{UPC repository},
Conjectures~4.1 and~4.5, pp.~51 and~60.
\end{thebibliography}
\end{document}
